\section{Discussion and conclusion}\label{sec:discussion}

\paragraph{Which choices mattered.} The number of particles and the decision to resample at all decided most of
the accuracy. The choices made inside resampling decided particle diversity instead, and the two are not the same
thing: a difference in diversity reaches the estimate only when particles are scarce and each one carries a map.
The proposal sits between them, mattering most where particles are few and losing its advantage as they are added.
A measurement-informed proposal also costs more per particle, so the fair comparison is at equal runtime rather
than at equal particle count, and there its advantage is much smaller. Adaptive resampling buys no accuracy in
these worlds but saves resampling steps at no cost.

\paragraph{What limits learning.} Two limits, of different kinds. The first is the inference: the objective is
estimated by the same weights that do the filtering, so a filter that has not found the robot reports which
setting lets it work rather than which setting generated the data, and a filter with too few particles reports
noise. Inference quality is a precondition for learning, not merely something learning improves. The second is the
data: the motion parameters are perfectly well defined and undetermined by 30 steps of this world, and become
determined only when the trajectory is long enough for their effect to show against the measurement noise. No
amount of computation removes that, and agreement between independent recordings separates the two cases cheaply. The one-sided shape of the likelihood also explains a common practice: since understating a noise level is
penalized far more heavily than overstating it, values set by hand are rationally kept large, as the settings
inherited from the upstream code are.

\paragraph{On reusing implementations.} Reused code has to be checked function by function before it is
compared: four errors in the upstream code would each have biased a comparison, one of them leaving resampled
particles sharing their state, which would have masked exactly the diversity effect measured here.

\paragraph{Limitations.} The worlds are small and simulated, with known correspondences; real sensors, unknown
data association and larger maps would make degeneracy and proposal quality more important. Each experiment
varies one choice while the others are fixed, so interactions are not measured. The rules for when to resample
were compared only in localization, and the resampling schemes in SLAM only with FastSLAM 1.0. The proposal and
the number of particles were held fixed rather than varied, so nothing here says how the two choices studied
interact with them. The learning
experiments estimate one parameter at a time from a known start, so they do not give the joint maximum over all
four, and a sweep is a cruder optimizer than expectation maximization.

\paragraph{Conclusion.} Localization and SLAM, written as a dynamic Bayesian network, make the same three demands
as any graphical model: a structure that justifies a factorization, an inference procedure for what the structure
leaves open, and parameters that have to come from somewhere. This project took the structure as given and worked
on the other two. On the inference side it measured which sampling choices change the answer, and found that most
of the accuracy is decided before the resampling scheme is chosen. On the learning side it estimated the model's
own noise parameters from recorded runs, recovered the measurement noise to within one step of the sweep, and
showed that the motion noise is identifiable only with a longer trajectory. The two sides meet in one result: the
quality of the approximate inference sets the ceiling on what can be learned.

With more time the obvious next steps would be MCMC moves after resampling \cite{gilks2001following}, which
attack particle impoverishment directly; the rules for when to resample in SLAM; estimating the four parameters
together by expectation maximization instead of one sweep at a time; and running the whole thing on real data,
where the correspondences are not given.
